\section{Introduction}
\IEEEPARstart{C}{onverter-rich} power systems couple controller parameters,
active-power schedules, and network conditions through both the ac operating
point and the dynamic Jacobian. The effect of changing two devices together
therefore need not equal the sum of changing them separately. This elementary
observation is operationally important: screening tools are frequently local
or one-action-at-a-time, while actual planning and tuning decisions arrive as
finite portfolios.

Several mature theories already illuminate parts of this problem. Classical
modal analysis relates states, residues, participations, and eigenvalue
sensitivities to oscillatory behavior
\cite{perezarriaga1982sma,pagola1989sensitivities}. Augmented and
multi-parameter sensitivities can retain equilibrium dependence and
higher-order terms \cite{nam2000eigensensitivity,zamoracardenas2013multiparameter}.
Return-ratio and impedance methods expose feedback interactions without
requiring the same internal coordinate description
\cite{macfarlane1970return,sun2011impedance,rygg2016sequence}. Recent work
also uses open-loop modal proximity and second-order eigenvalue displacement to
identify converter interaction modes \cite{dewangan2026interaction}.

These methods answer valuable but different questions. A derivative is local
to its expansion point. An impedance or return-ratio criterion evaluates a
feedback interconnection. A participation factor associates a mode with
states. None of these objects, by itself, is the finite inclusion--exclusion
remainder of a specified engineering portfolio after every subset has been
physically applied and independently re-equilibrated. The combinatorial
operation needed to isolate that remainder is M\"obius inversion
\cite{rota1964mobius,grabisch1997kadditive}. The missing bridge is to apply it
to a re-equilibrated differential--algebraic operator and then connect the
finite result to derivative, factor, and modal descriptions without confusing
existence with consequence.

This paper develops and tests that bridge. A normalized action coordinate
maps zero to a registered baseline and one to a finite engineering change.
For each coalition vertex, the ac power flow is solved again, dynamic states
are initialized again, and the descriptor operator is rebuilt. A relative
real log-determinant, integrated over a fixed frequency band, summarizes the
resulting operator. Direct pair and triple M\"obius differences define the
finite externality. Independently, a connected set-partition formula
integrates mixed total derivatives over the action hypercube. The two routes
must close within a case-specific uncertainty budget.

The contribution is not a new M\"obius identity, Schur complement, trace
formula, perturbation determinant, or eigenvalue sensitivity. It is the
following integrated, experimentally gated workflow:
\begin{enumerate}
    \item a finite set function whose every vertex is an independently
    re-equilibrated power-system DAE;
    \item a connected reconstruction that retains mixed physical and
    equilibrium curvature rather than only singleton feedback;
    \item an exact descriptor-factor anatomy separating algebraic,
    mass-scaling, and finite-pole externalities;
    \item a registered modal-materiality ladder that asks whether the
    externality reaches the mode that actually sets the damping margin.
\end{enumerate}

Table~\ref{tab:methods} makes the novelty perimeter explicit. The first four
rows are established approaches; none is displaced here. M\"obius inversion
supplies the set-function coefficient, higher-order sensitivity supplies local
derivative information, and return-ratio/impedance methods supply loop-level
stability information. The proposed contribution is the registered
combination of a finite re-equilibrated power-system set function, complete
pair/triple subset evaluation, connected reconstruction, descriptor-factor
anatomy, and a separate materiality test.

\begin{table*}[t]
\caption{Method-to-Method Novelty Perimeter}
\label{tab:methods}
\centering
\scriptsize
\begin{tabular}{p{0.17\textwidth}p{0.12\textwidth}p{0.16\textwidth}p{0.21\textwidth}p{0.19\textwidth}}
\toprule
Approach & Amplitude & Equilibrium treatment & Interaction object & Primary interpretation\\
\midrule
First-order eigen-sensitivity
& Local
& One expansion point
& Singleton derivative of a selected eigenvalue
& Directional modal change\\
Higher-order total sensitivity
& Local Taylor series
& May include implicit equilibrium derivatives
& Selected mixed derivatives to a truncation order
& Local curvature and parameter ranking\\
Return-ratio / impedance
& Frequency-domain model
& Operating point supplied to the model
& Ports, loops, or interconnection blocks
& Nyquist/loop stability and coupling\\
Proximal eigenvalue displacement \cite{dewangan2026interaction}
& First-/second-order modal displacement
& Not an all-subset finite-action requirement
& Proximal open-loop modal coupling
& Identification of converter interaction modes\\
M\"obius set-function coefficient
& Finite
& Domain-agnostic vertex values
& All subsets of a specified set
& Pure higher-order set-function dividend\\
This work
& Finite registered actions
& New AC solution and DAE at every required vertex
& All vertices of every registered pair and triple
& Externality, exact factor anatomy, then separate registered modal materiality\\
\bottomrule
\end{tabular}
\end{table*}

The tests use a canonical IEEE 39-bus system with three grid-following (GFL)
converters and eight moderate actions. The main positive result is numerical
and structural: reproducible pair and genuine third-order finite-pole
externalities exist, and the connected reconstruction closes the direct
finite computation. The main negative result is equally important: the
largest externalities found in the registered search do not materially alter
the limiting damping mode. A post-mortem identifies a family mismatch. The
externality-dominant modes are well-damped converter-control or
phase-locked-loop (PLL) modes, while the limiting mode is a lightly damped
synchronous-electromechanical mode. Hence \emph{dynamic externality existence
and operational interpretation under a registered modal-materiality screen
are distinct claims}.

This distinction narrows the practical interpretation. The calculus can tell
an analyst that a finite portfolio contains reproducible, non-additive dynamic
structure, quantify what a frozen model misses, and locate pole motion when a
contour certificate is valid. It does not certify instability, identify an
industrial mitigation, establish cycle or strongly connected component
(SCC) causality, or transfer a mechanism to electromagnetic-transient (EMT)
simulation. The cross-domain EMT work in this study validates only the matched
benchmark infrastructure used to launch the phasor-domain campaign.

The remainder of the paper is organized as follows. Section II defines the
finite calculus and exact factorization. Section III describes the registered
benchmark, actions, uncertainty budgets, and holdout discipline. Section IV
reports finite externalities and connected reconstruction. Section V separates
algebraic and pole contributions and presents contour localization. Section VI
reports the modal consequence and materiality boundary. Sections VII and VIII
discuss interpretation, limitations, and conclusions.
